\chapter*{Preface}
\addcontentsline{toc}{chapter}{Preface}
\markboth{Preface}{Preface}

A home-grown protocol is an instructive thing to build. You decide for yourself where the frame
begins, how long it is, how it is protected against errors, and what should happen when a reply
never comes. Every mechanism becomes your own --- and every mechanism becomes something that has
to be written and maintained.

This book is about what happens when someone else has already decided all of that, once, and cast
it in silicon.

CAN dates from 1986, it is in every car built since the mid-nineties, and it solves almost exactly
the problems you have just solved yourselves --- only in hardware, and better. Framing, checksum,
collision handling and acknowledgement all live in the controller. What is left for the software
is to hand over an identifier, a length and a few data bytes, and then read back how it went.

\begin{aside}
\note That CAN does this for you does not mean you get away without understanding it. If you build
the controller, the protocol is your requirements specification, bit for bit. If you write the
driver, you have to handle precisely what the controller \emph{cannot} tell you: why a lost
arbitration looks like a successful transmission, why a frame that nobody acknowledges completes
anyway, and why a received frame can vanish from under your nose. All of that follows from how CAN
works, and it is all in this book.
\end{aside}

\section*{What the book covers}

\Kapref{ch:buss} starts at the wires: why a shared bus, what dominant and recessive bits are, and
why that asymmetry is the whole foundation. \Kapref{ch:frameformat} goes through the CAN frame
field by field. \Kapref{ch:stuffing} covers bit stuffing and the CRC-15, the two mechanisms that
let a frame survive a noisy bus. \Kapref{ch:arbitrering} works through arbitration bit by bit, by
hand. \Kapref{ch:tajming} is about bit timing and the sample point: how nodes without a shared
clock still read the same bit at the same time. \Kapref{ch:fel} describes error handling in real
CAN --- error frames, error counters and bus-off.

\Kapref{ch:kursen} is the only chapter in the book that is not about CAN in general. It describes
the simplified controller that is built in teaching, what it implements of all that comes before,
and what it deliberately leaves out. That chapter is the one that comes in useful when debugging
--- from either side.

\section*{What the book does not cover}

The book describes \textbf{classical CAN} with 11-bit identifiers, that is CAN 2.0A. Extended
29-bit identifiers (2.0B) are mentioned where they affect the shape of the frame, but are not
worked through. CAN~FD and CAN~XL are out of scope, as are higher layers such as CANopen and
J1939.

Nor does the book describe the electrical side in detail: differential signalling, the physics of
termination and the effect of cable length are touched on only as far as is needed to understand
why the bus looks the way it does.

\section*{Exercises}

Every chapter ends with a few exercises. They are done with pen and paper --- there is no code to
compile in this book --- and the answers to all of them are in Appendix~A.

The exercises are short on purpose. They are there to check that the chapter stuck, not to take up
an evening.
